\IfFileExists{stacks-project.cls}{%
\documentclass{stacks-project}
}{%
\documentclass{amsart}
}

% For dealing with references we use the comment environment
\usepackage{verbatim}
\newenvironment{reference}{\comment}{\endcomment}
%\newenvironment{reference}{}{}
\newenvironment{slogan}{\comment}{\endcomment}
\newenvironment{history}{\comment}{\endcomment}

% For commutative diagrams we use Xy-pic
\usepackage[all]{xy}

% We use 2cell for 2-commutative diagrams.
\xyoption{2cell}
\UseAllTwocells

% We use multicol for the list of chapters between chapters
\usepackage{multicol}

% This is generally recommended for better output
\usepackage{lmodern}
\usepackage[T1]{fontenc}

% For cross-file-references
\usepackage{xr-hyper}

% Package for hypertext links:
\usepackage{hyperref}

% For any local file, say "hello.tex" you want to link to please
% use \externaldocument[hello-]{hello}
\externaldocument[introduction-]{introduction}
\externaldocument[conventions-]{conventions}
\externaldocument[sets-]{sets}
\externaldocument[categories-]{categories}
\externaldocument[topology-]{topology}
\externaldocument[sheaves-]{sheaves}
\externaldocument[sites-]{sites}
\externaldocument[stacks-]{stacks}
\externaldocument[fields-]{fields}
\externaldocument[algebra-]{algebra}
\externaldocument[brauer-]{brauer}
\externaldocument[homology-]{homology}
\externaldocument[derived-]{derived}
\externaldocument[simplicial-]{simplicial}
\externaldocument[more-algebra-]{more-algebra}
\externaldocument[smoothing-]{smoothing}
\externaldocument[modules-]{modules}
\externaldocument[sites-modules-]{sites-modules}
\externaldocument[injectives-]{injectives}
\externaldocument[cohomology-]{cohomology}
\externaldocument[sites-cohomology-]{sites-cohomology}
\externaldocument[dga-]{dga}
\externaldocument[dpa-]{dpa}
\externaldocument[sdga-]{sdga}
\externaldocument[hypercovering-]{hypercovering}
\externaldocument[schemes-]{schemes}
\externaldocument[constructions-]{constructions}
\externaldocument[properties-]{properties}
\externaldocument[morphisms-]{morphisms}
\externaldocument[coherent-]{coherent}
\externaldocument[divisors-]{divisors}
\externaldocument[limits-]{limits}
\externaldocument[varieties-]{varieties}
\externaldocument[topologies-]{topologies}
\externaldocument[descent-]{descent}
\externaldocument[perfect-]{perfect}
\externaldocument[more-morphisms-]{more-morphisms}
\externaldocument[flat-]{flat}
\externaldocument[groupoids-]{groupoids}
\externaldocument[more-groupoids-]{more-groupoids}
\externaldocument[etale-]{etale}
\externaldocument[chow-]{chow}
\externaldocument[intersection-]{intersection}
\externaldocument[pic-]{pic}
\externaldocument[weil-]{weil}
\externaldocument[adequate-]{adequate}
\externaldocument[dualizing-]{dualizing}
\externaldocument[duality-]{duality}
\externaldocument[discriminant-]{discriminant}
\externaldocument[derham-]{derham}
\externaldocument[local-cohomology-]{local-cohomology}
\externaldocument[algebraization-]{algebraization}
\externaldocument[curves-]{curves}
\externaldocument[resolve-]{resolve}
\externaldocument[models-]{models}
\externaldocument[functors-]{functors}
\externaldocument[equiv-]{equiv}
\externaldocument[pione-]{pione}
\externaldocument[etale-cohomology-]{etale-cohomology}
\externaldocument[proetale-]{proetale}
\externaldocument[relative-cycles-]{relative-cycles}
\externaldocument[more-etale-]{more-etale}
\externaldocument[trace-]{trace}
\externaldocument[crystalline-]{crystalline}
\externaldocument[spaces-]{spaces}
\externaldocument[spaces-properties-]{spaces-properties}
\externaldocument[spaces-morphisms-]{spaces-morphisms}
\externaldocument[decent-spaces-]{decent-spaces}
\externaldocument[spaces-cohomology-]{spaces-cohomology}
\externaldocument[spaces-limits-]{spaces-limits}
\externaldocument[spaces-divisors-]{spaces-divisors}
\externaldocument[spaces-over-fields-]{spaces-over-fields}
\externaldocument[spaces-topologies-]{spaces-topologies}
\externaldocument[spaces-descent-]{spaces-descent}
\externaldocument[spaces-perfect-]{spaces-perfect}
\externaldocument[spaces-more-morphisms-]{spaces-more-morphisms}
\externaldocument[spaces-flat-]{spaces-flat}
\externaldocument[spaces-groupoids-]{spaces-groupoids}
\externaldocument[spaces-more-groupoids-]{spaces-more-groupoids}
\externaldocument[bootstrap-]{bootstrap}
\externaldocument[spaces-pushouts-]{spaces-pushouts}
\externaldocument[spaces-chow-]{spaces-chow}
\externaldocument[groupoids-quotients-]{groupoids-quotients}
\externaldocument[spaces-more-cohomology-]{spaces-more-cohomology}
\externaldocument[spaces-simplicial-]{spaces-simplicial}
\externaldocument[spaces-duality-]{spaces-duality}
\externaldocument[formal-spaces-]{formal-spaces}
\externaldocument[restricted-]{restricted}
\externaldocument[spaces-resolve-]{spaces-resolve}
\externaldocument[formal-defos-]{formal-defos}
\externaldocument[defos-]{defos}
\externaldocument[cotangent-]{cotangent}
\externaldocument[examples-defos-]{examples-defos}
\externaldocument[algebraic-]{algebraic}
\externaldocument[examples-stacks-]{examples-stacks}
\externaldocument[stacks-sheaves-]{stacks-sheaves}
\externaldocument[criteria-]{criteria}
\externaldocument[artin-]{artin}
\externaldocument[quot-]{quot}
\externaldocument[stacks-properties-]{stacks-properties}
\externaldocument[stacks-morphisms-]{stacks-morphisms}
\externaldocument[stacks-limits-]{stacks-limits}
\externaldocument[stacks-cohomology-]{stacks-cohomology}
\externaldocument[stacks-perfect-]{stacks-perfect}
\externaldocument[stacks-introduction-]{stacks-introduction}
\externaldocument[stacks-more-morphisms-]{stacks-more-morphisms}
\externaldocument[stacks-geometry-]{stacks-geometry}
\externaldocument[moduli-]{moduli}
\externaldocument[moduli-curves-]{moduli-curves}
\externaldocument[examples-]{examples}
\externaldocument[exercises-]{exercises}
\externaldocument[guide-]{guide}
\externaldocument[desirables-]{desirables}
\externaldocument[coding-]{coding}
\externaldocument[obsolete-]{obsolete}
\externaldocument[fdl-]{fdl}
\externaldocument[index-]{index}

% Theorem environments.
%
\theoremstyle{plain}
\newtheorem{theorem}[subsection]{Theorem}
\newtheorem{proposition}[subsection]{Proposition}
\newtheorem{lemma}[subsection]{Lemma}

\theoremstyle{definition}
\newtheorem{definition}[subsection]{Definition}
\newtheorem{example}[subsection]{Example}
\newtheorem{exercise}[subsection]{Exercise}
\newtheorem{situation}[subsection]{Situation}

\theoremstyle{remark}
\newtheorem{remark}[subsection]{Remark}
\newtheorem{remarks}[subsection]{Remarks}

\numberwithin{equation}{subsection}

% Macros
%
\def\lim{\mathop{\mathrm{lim}}\nolimits}
\def\colim{\mathop{\mathrm{colim}}\nolimits}
\def\Spec{\mathop{\mathrm{Spec}}}
\def\Hom{\mathop{\mathrm{Hom}}\nolimits}
\def\Ext{\mathop{\mathrm{Ext}}\nolimits}
\def\SheafHom{\mathop{\mathcal{H}\!\mathit{om}}\nolimits}
\def\SheafExt{\mathop{\mathcal{E}\!\mathit{xt}}\nolimits}
\def\Sch{\mathit{Sch}}
\def\Mor{\mathop{\mathrm{Mor}}\nolimits}
\def\Ob{\mathop{\mathrm{Ob}}\nolimits}
\def\Sh{\mathop{\mathit{Sh}}\nolimits}
\def\NL{\mathop{N\!L}\nolimits}
\def\CH{\mathop{\mathrm{CH}}\nolimits}
\def\proetale{{pro\text{-}\acute{e}tale}}
\def\etale{{\acute{e}tale}}
\def\QCoh{\mathit{QCoh}}
\def\Ker{\mathop{\mathrm{Ker}}}
\def\Im{\mathop{\mathrm{Im}}}
\def\Coker{\mathop{\mathrm{Coker}}}
\def\Coim{\mathop{\mathrm{Coim}}}

% Boxtimes
%
\DeclareMathSymbol{\boxtimes}{\mathbin}{AMSa}{"02}

%
% Macros for moduli stacks/spaces
%
\def\QCohstack{\mathcal{QC}\!\mathit{oh}}
\def\Cohstack{\mathcal{C}\!\mathit{oh}}
\def\Spacesstack{\mathcal{S}\!\mathit{paces}}
\def\Quotfunctor{\mathrm{Quot}}
\def\Hilbfunctor{\mathrm{Hilb}}
\def\Curvesstack{\mathcal{C}\!\mathit{urves}}
\def\Polarizedstack{\mathcal{P}\!\mathit{olarized}}
\def\Complexesstack{\mathcal{C}\!\mathit{omplexes}}
% \Pic is the operator that assigns to X its picard group, usage \Pic(X)
% \Picardstack_{X/B} denotes the Picard stack of X over B
% \Picardfunctor_{X/B} denotes the Picard functor of X over B
\def\Pic{\mathop{\mathrm{Pic}}\nolimits}
\def\Picardstack{\mathcal{P}\!\mathit{ic}}
\def\Picardfunctor{\mathrm{Pic}}
\def\Deformationcategory{\mathcal{D}\!\mathit{ef}}
\usepackage{cleveref}
\crefname{originaldefinition}{original definition}{original definitions}
\Crefname{originaldefinition}{Original definition}{Original definitions}
\makeatletter
\newlabel{definition-reduced-induced-scheme}{{01J4}{source}{}{sourcepointer-1}{}}
\newlabel{definition-reduced-induced-scheme@cref}{{[originaldefinition][0][]01J4}{[1][0][]source}{}{}{}}
\makeatother
\crefname{originallemma}{original lemma}{original lemmas}
\Crefname{originallemma}{Original lemma}{Original lemmas}
\makeatletter
\newlabel{lemma-image-quasi-compact-closed}{{05JL}{source}{}{sourcepointer-2}{}}
\newlabel{lemma-image-quasi-compact-closed@cref}{{[originallemma][0][]05JL}{[1][0][]source}{}{}{}}
\makeatother
\crefname{originalreference}{original reference}{original references}
\Crefname{originalreference}{Original reference}{Original references}
\makeatletter
\newlabel{topology-definition-lift-specializations}{{0063}{source}{}{sourcepointer-3}{}}
\newlabel{topology-definition-lift-specializations@cref}{{[originalreference][0][]0063}{[1][0][]source}{}{}{}}
\makeatother
\crefname{originalreference}{original reference}{original references}
\Crefname{originalreference}{Original reference}{Original references}
\makeatletter
\newlabel{topology-lemma-closed-open-map-specialization}{{0066}{source}{}{sourcepointer-4}{}}
\newlabel{topology-lemma-closed-open-map-specialization@cref}{{[originalreference][0][]0066}{[1][0][]source}{}{}{}}
\makeatother
\crefname{originalreference}{original reference}{original references}
\Crefname{originalreference}{Original reference}{Original references}
\makeatletter
\newlabel{topology-lemma-lift-specialization-composition}{{0064}{source}{}{sourcepointer-5}{}}
\newlabel{topology-lemma-lift-specialization-composition@cref}{{[originalreference][0][]0064}{[1][0][]source}{}{}{}}
\makeatother
\crefname{originalreference}{original reference}{original references}
\Crefname{originalreference}{Original reference}{Original references}
\makeatletter
\newlabel{topology-lemma-lift-specializations-images}{{0065}{source}{}{sourcepointer-6}{}}
\newlabel{topology-lemma-lift-specializations-images@cref}{{[originalreference][0][]0065}{[1][0][]source}{}{}{}}
\makeatother
\title{For qc morphisms, closedness is the same as specialization lifting}
\author{The Stacks Project Authors}
\date{stacks-project:schemes.tex:local-snapshot-20261012; GNU FDL 1.2 or later}
\begin{document}
\maketitle
\noindent\textit{Editorial note.} Counted span is the author's own proof environment, source lines 3916-3933, following the statement at 3908-3914.\par\medskip
\noindent\textit{Editorial note.} The step 'the restriction of f to Z is still quasi-compact' (\textasciitilde{}3924-3925) is uncited but standard (a closed immersion is quasi-compact and qc is preserved under composition).\par\medskip
\noindent\textit{Allowed standard prerequisites (editorial).}\begin{itemize}
\item Closed maps lift specializations. Topology, Lemma topology-lemma-closed-open-map-specialization (cross-chapter, cited at \textasciitilde{}3918). (1) direction.
\item Reduced induced scheme structure. definition-reduced-induced-scheme (same chapter, cited at \textasciitilde{}3923). Views the closed subset Z as a scheme.
\item Specialization lifting and composition. Topology, Lemma topology-lemma-lift-specialization-composition (cross-chapter, cited at \textasciitilde{}3927). Specializations lift along Z -> S.
\item Images of qc maps stable under specializations. Topology, Lemma topology-lemma-lift-specializations-images (cross-chapter, cited at \textasciitilde{}3930). f(X) is specialization-stable.
\item Specialization-stable image of a qc map is closed. lemma-image-quasi-compact-closed (same chapter, cited at \textasciitilde{}3932). Concluding engine.
\end{itemize}\medskip
\ifdefined\hypertarget\hypertarget{sourcepointer-1}{}\fi
\noindent\textit{Source reference.} Original definition 01J4: sources/stacks-schemes/pinned-original-source.tex, source lines 2456--2456. Navigation pointer only: the counted proof cites this statement by \textbackslash{}ref as a standing input; the deduction the counted proof has to make is carried out inside the counted span, and no part of that deduction is delegated to this pointer.\par\medskip
\ifdefined\hypertarget\hypertarget{sourcepointer-2}{}\fi
\noindent\textit{Source reference.} Original lemma 05JL: sources/stacks-schemes/pinned-original-source.tex, source lines 3883--3883. Navigation pointer only: the counted proof cites this statement by \textbackslash{}ref as a standing input; the deduction the counted proof has to make is carried out inside the counted span, and no part of that deduction is delegated to this pointer.\par\medskip
\ifdefined\hypertarget\hypertarget{sourcepointer-3}{}\fi
\noindent\textit{Source reference.} Original reference 0063: focused-seeds/geometry/S9-schemes-quasi-compact-closed/ref-tags.txt, source lines 5--5. Navigation pointer only: the counted proof cites this statement by \textbackslash{}ref as a standing input; the deduction the counted proof has to make is carried out inside the counted span, and no part of that deduction is delegated to this pointer.\par\medskip
\ifdefined\hypertarget\hypertarget{sourcepointer-4}{}\fi
\noindent\textit{Source reference.} Original reference 0066: focused-seeds/geometry/S9-schemes-quasi-compact-closed/ref-tags.txt, source lines 6--6. Navigation pointer only: the counted proof cites this statement by \textbackslash{}ref as a standing input; the deduction the counted proof has to make is carried out inside the counted span, and no part of that deduction is delegated to this pointer.\par\medskip
\ifdefined\hypertarget\hypertarget{sourcepointer-5}{}\fi
\noindent\textit{Source reference.} Original reference 0064: focused-seeds/geometry/S9-schemes-quasi-compact-closed/ref-tags.txt, source lines 7--7. Navigation pointer only: the counted proof cites this statement by \textbackslash{}ref as a standing input; the deduction the counted proof has to make is carried out inside the counted span, and no part of that deduction is delegated to this pointer.\par\medskip
\ifdefined\hypertarget\hypertarget{sourcepointer-6}{}\fi
\noindent\textit{Source reference.} Original reference 0065: focused-seeds/geometry/S9-schemes-quasi-compact-closed/ref-tags.txt, source lines 8--8. Navigation pointer only: the counted proof cites this statement by \textbackslash{}ref as a standing input; the deduction the counted proof has to make is carried out inside the counted span, and no part of that deduction is delegated to this pointer.\par\medskip
\begin{lemma}
\label{lemma-quasi-compact-closed}
Let $f : X \to S$ be a quasi-compact morphism of schemes.
Then $f$ is closed if and only if specializations lift
along $f$, see
Topology, Definition \ref{topology-definition-lift-specializations}.
\end{lemma}

\begin{proof}
According to
Topology, Lemma \ref{topology-lemma-closed-open-map-specialization}
if $f$ is closed then specializations lift along $f$.
Conversely, suppose that specializations lift along $f$.
Let $Z \subset X$ be a closed subset. We may think of $Z$
as a scheme with the reduced induced scheme structure, see
Definition \ref{definition-reduced-induced-scheme}.
Since $Z \subset X$ is closed the restriction
of $f$ to $Z$ is still quasi-compact. Moreover specializations lift
along $Z \to S$ as well,
see Topology, Lemma \ref{topology-lemma-lift-specialization-composition}.
Hence it suffices to prove $f(X)$ is closed if specializations lift along $f$.
In particular $f(X)$ is stable under specializations, see
Topology, Lemma \ref{topology-lemma-lift-specializations-images}.
Thus $f(X)$ is closed by
Lemma \ref{lemma-image-quasi-compact-closed}.
\end{proof}

\end{document}
